% =====================================================================================
% La puerta de Hadamard vista sobre la esfera de Bloch — figura propia en TikZ.
%
% Tres viñetas: el estado de partida (|0>+|1>)/raiz(2) sobre el eje x, la rotacion de
% 90 grados en torno a y que lo lleva al polo sur, y la rotacion de 180 grados en torno
% a x que lo sube al polo norte. Resultado: H|+> = |0>.
%
% CODIGO DE COLOR, deliberado y constante en toda la memoria:
%   azulunir   -> el QUBIT (vector de estado, punto y etiqueta)
%   acentofig  -> las TRANSFORMACIONES (giros, planos, flechas de operacion)
%
% ⚠️ \shorthandoff{<>} es obligatorio: `babel` en español hace activos < y >, y eso
%    rompe la sintaxis de puntas de flecha de TikZ.
% =====================================================================================
\begingroup
\shorthandoff{<>}%
\begin{tikzpicture}[>={Stealth[length=2mm]}, line width=0.5pt, scale=0.92]

  \pgfmathsetmacro{\R}{1.45}    % radio de cada esfera
  \pgfmathsetmacro{\K}{0.30}    % achatamiento del ecuador; ver nota al pie del fichero
  \pgfmathsetmacro{\Tx}{235}    % angulo de pantalla del eje x
  \pgfmathsetmacro{\D}{4.9}     % separacion entre viñetas

  % --- una esfera con sus tres ejes, reutilizable ------------------------------------
  \newcommand{\esferabase}{%
    \draw (0,0) circle (\R);
    \draw[dashed] (0,0) ellipse ({\R} and {\K*\R});
    \draw[->] (0,0) -- (0,\R);
    \node[right=1pt] at ({0.02*\R},{0.78*\R}) {\scriptsize $\hat{\mathbf{z}}$};
    \draw[->] (0,0) -- (\R,0) node[right=-1pt] {\scriptsize $\hat{\mathbf{y}}$};
    \draw[->] (0,0) -- ({\R*cos(\Tx)},{\K*\R*sin(\Tx)})
          node[below left=-3pt] {\scriptsize $\hat{\mathbf{x}}$};%
  }

  % =====================================================================================
  % 1) Estado de partida: (|0>+|1>)/raiz(2), sobre el eje x
  % =====================================================================================
  \begin{scope}[xshift=0cm]
    \esferabase
    \draw[azulunir, line width=1.1pt]
          (0,0) -- ({\R*cos(\Tx)},{\K*\R*sin(\Tx)});
    \fill[azulunir] ({\R*cos(\Tx)},{\K*\R*sin(\Tx)}) circle (2pt);
    \node[azulunir, anchor=east] at ({-1.06*\R},{-0.40*\R})
         {\small $(\ket{0}+\ket{1})/\sqrt{2}$};
    % giro de 90 grados en torno al eje y
    \draw[acentofig, ->, line width=0.7pt]
          ({0.30*\R},{-0.038*\R}) arc[start angle=-170, end angle=110,
          x radius={0.14*\R}, y radius={0.22*\R}];
  \end{scope}

  \draw[->, line width=0.8pt] ({0.5*\D-0.45},0) -- ({0.5*\D+0.45},0);

  % =====================================================================================
  % 2) Tras el giro: |1> en el polo sur. La rotacion de 180 grados en torno a x se
  %    representa con el plano y la flecha doble.
  % =====================================================================================
  \begin{scope}[xshift=\D cm]
    % el plano va PRIMERO: asi la esfera y el estado quedan por encima y se ven
    \draw[acentofig, fill=acentofig, fill opacity=0.06, line width=0.6pt]
          ({-1.78*\R},{-0.33*\R}) -- ({-0.90*\R},{0.31*\R})
       -- ({ 1.78*\R},{ 0.31*\R}) -- ({ 0.90*\R},{-0.33*\R}) -- cycle;
    \esferabase
    \draw[azulunir, line width=1.1pt] (0,0) -- (0,-\R);
    \fill[azulunir] (0,-\R) circle (2pt);
    \node[azulunir, anchor=north] at (0,{-1.08*\R}) {\small $\ket{1}$};
    % flecha doble: el estado atraviesa el plano de arriba abajo
    \draw[acentofig, <->, line width=0.9pt]
          ({0.58*\R},{-0.52*\R}) -- ({0.58*\R},{0.52*\R});
  \end{scope}

  \draw[->, line width=0.8pt] ({1.5*\D-0.45},0) -- ({1.5*\D+0.45},0);

  % =====================================================================================
  % 3) Resultado: |0> en el polo norte
  % =====================================================================================
  \pgfmathsetmacro{\DD}{2*\D}
    \begin{scope}[xshift=\DD cm]
    \esferabase
    \draw[azulunir, line width=1.1pt] (0,0) -- (0,\R);
    \fill[azulunir] (0,\R) circle (2pt);
    \node[azulunir, anchor=south east] at ({-0.04*\R},{1.06*\R}) {\small $\ket{0}$};
  \end{scope}

\end{tikzpicture}%
\endgroup

% =====================================================================================
% NOTA SOBRE \K, EL ACHATAMIENTO
%
% Determina desde que altura se mira la esfera, y con ello el angulo VISUAL entre los
% ejes x e y. Con el eje y horizontal, ese angulo no puede llegar a 90 grados: el eje x
% ha de apoyarse sobre la elipse del ecuador —porque el ecuador ES el plano xy— y la
% unica forma de que forme 90 grados con y seria que apuntase recto hacia abajo, encima
% del eje z. Lo medido:
%
%     K = 0,30  ->  157 grados   punta de x a 0,62·R del centro
%     K = 0,55  ->  142 grados   punta de x a 0,73·R del centro   <- el elegido
%     K = 0,75  ->  133 grados   punta de x a 0,84·R del centro
%     K = 1,00  ->  125 grados   punta de x a 1,00·R  (el ecuador se confunde con
%                                el contorno y la esfera deja de parecer una esfera)
%
% Se probo 0,55 y se descarto: el usuario prefiere la perspectiva original de 0,30.
% La tabla se deja aqui por si en algun momento se quiere volver a mover.
% =====================================================================================
